\section{Email Templates}

\begin{frame}[fragile]{What Is an Email Template?}
	\begin{itemize}
		\item An \textbf{email template} is a reusable mail design stored in \texttt{mail.template}.
		\item It can render subject/body with QWeb and record values.
		\item Used for welcome mails, resets, invoices, reminders, and notifications.
\begin{block}{Model}
\begin{lstlisting}
mail.template
\end{lstlisting}
\end{block}
	\end{itemize}
\end{frame}

\begin{frame}[fragile]{Defining an Email Template in XML}
\begin{lstlisting}[style=xml]
<record id="mail_template_student_welcome" model="mail.template">
  <field name="name">Student: Welcome Email</field>
  <field name="model_id" ref="model_student_student"/>
  <field name="subject">Welcome ${object.name}</field>
  <field name="email_from">{{ (object.company_id.email or user.email) }}</field>
  <field name="email_to">{{ object.email }}</field>
  <field name="body_html" type="html">
	<p>Hello <t t-out="object.name"/>,</p>
	<p>Welcome to our school system.</p>
  </field>
</record>
\end{lstlisting}
\end{frame}

\begin{frame}[fragile]{Breaking Down Template Fields}
	\begin{itemize}
		\item \textbf{\texttt{name}}: template label
		\item \textbf{\texttt{model\_id}}: model this template applies to
		\item \textbf{\texttt{subject}}: email subject (dynamic)
		\item \textbf{\texttt{email\_from}} / \textbf{\texttt{email\_to}}: addresses
		\item \textbf{\texttt{body\_html}}: HTML body, often QWeb
		\item Optional: \texttt{report\_template\_ids}, attachments, language, etc.
	\end{itemize}
\end{frame}

\begin{frame}[fragile]{Sending a Template from Python}
\begin{block}{Method Syntax}
\begin{lstlisting}
def action_send_welcome_email(self):
	template = self.env.ref(
		'school_manager.mail_template_student_welcome'
	)
	for student in self:
		template.send_mail(student.id, force_send=True)
	return True
\end{lstlisting}
\end{block}
	\begin{itemize}
		\item \texttt{send\_mail(res\_id)} renders and sends for that record.
		\item \texttt{force\_send=True} sends immediately instead of only queueing.
	\end{itemize}
\end{frame}

\begin{frame}[fragile]{QWeb in Email Bodies}
\begin{lstlisting}[style=xml]
<field name="body_html" type="html">
  <div>
	<p>Dear <t t-out="object.name"/>,</p>
	<p>Your code is
	  <strong><t t-out="object.code"/></strong>.
	</p>
	<p t-if="object.classroom_id">
	  Classroom:
	  <t t-out="object.classroom_id.name"/>
	</p>
  </div>
</field>
\end{lstlisting}
	\begin{itemize}
		\item \texttt{object} is the current record.
		\item Use \texttt{t-out} in modern Odoo QWeb.
	\end{itemize}
\end{frame}

\begin{frame}[fragile]{Opening Composer with a Template}
\begin{lstlisting}
def action_share_by_email(self):
	self.ensure_one()
	template = self.env.ref(
		'school_manager.mail_template_student_welcome'
	)
	return {
		'type': 'ir.actions.act_window',
		'name': 'Send Email',
		'res_model': 'mail.compose.message',
		'view_mode': 'form',
		'target': 'new',
		'context': {
			'default_model': 'student.student',
			'default_res_ids': self.ids,
			'default_template_id': template.id,
			'default_composition_mode': 'comment',
		},
	}
\end{lstlisting}
\end{frame}

\begin{frame}{Email Template Best Practices}
	\begin{itemize}
		\item Keep one template per clear business event.
		\item Use XML IDs and \texttt{env.ref()} instead of hard-coded IDs.
		\item Prefer QWeb dynamic fields over hardcoded names.
		\item Test rendering with real demo records.
		\item Put templates under \texttt{data/} and load them from the manifest.
	\end{itemize}
\end{frame}
